\documentclass[aspectratio=169]{beamer}
\usetheme{Jacquenetta}

\usepackage[utf8]{inputenc}
\usepackage{amsmath,amssymb}
\usepackage{booktabs}
% Save code blocks as boxes to avoid fragile-frame conflicts
\newsavebox{\codeboxa}
\newsavebox{\codeboxb}

% === Metadata ===================================================
\title{Jacquenetta}
\subtitle{A modern Beamer theme for scientific presentations}
\author{Samuel Manchajm}
\institute{GitHub --- Jacquenetta-beamer-theme}
\date{\today}

% === Logo test (colorbox, no external file needed) =============
\logo{%
  \colorbox{jqdark}{%
    \textcolor{white}{%
      \fontsize{6}{7}\selectfont\bfseries\quad MY INSTITUTION\quad%
    }%
  }%
}

\begin{document}

% ============================================================
% TITLE
% ============================================================
\titleframe
\logo{}  % <-- retire le logo du footer après le titre
          % Supprime cette ligne pour le garder sur tous les slides

% ============================================================
\section{Getting Started}
% ============================================================

\begin{frame}{Minimal example}
    Use the Jacquenetta theme with a single line in your preamble:

    \vspace{0.1cm}
    {\small\ttfamily
    \textbackslash documentclass[aspectratio=169]\{beamer\}\\
    \textbackslash usetheme\{Jacquenetta\}\\[2pt]
    \textbackslash title\{My Presentation\}\\
    \textbackslash author\{Jane Doe\}\\
    \textbackslash institute\{University of Somewhere\}\\[2pt]
    \textbackslash begin\{document\}\\
    \quad\textbackslash titleframe\\
    \quad\textbackslash begin\{frame\}\{Hello, world!\}\\
    \quad\quad Your content here.\\
    \quad\textbackslash end\{frame\}\\
    \textbackslash end\{document\}
    }
\end{frame}

\begin{frame}{Theme options}
    Jacquenetta supports the following options:

    \vspace{0.3cm}
    \begin{tabular}{ll}
        \toprule
        \textbf{Option} & \textbf{Effect} \\
        \midrule
        \texttt{accent=<color>} & Change the accent color \\
        \texttt{noborder}       & Disable the signature border \\
        \bottomrule
    \end{tabular}

    \vspace{0.5cm}
    {\small\ttfamily
    \textbackslash usetheme[accent=teal]\{Jacquenetta\}\\
    \textbackslash usetheme[noborder]\{Jacquenetta\}\\
    \textbackslash usetheme[accent=jqorange, noborder]\{Jacquenetta\}
    }
\end{frame}

% ============================================================
\section{Typography}
% ============================================================

\sectionframe{01}{Typography}

\begin{frame}{Font styles}
    Jacquenetta uses \textbf{Helvetica} for body text and
    \textit{serif fonts} for mathematics.

    \vspace{0.4cm}
    \begin{itemize}
        \item Regular text for general content
        \item \textbf{Bold text} for emphasis and headings
        \item \textit{Italic text} for definitions or foreign terms
        \item \texttt{Monospace} for code and technical identifiers
        \item \alert{Alerted text} for drawing attention
    \end{itemize}

    \vspace{0.4cm}
    {\footnotesize Footnote-sized text is also available for captions
    and annotations.}
\end{frame}

\begin{frame}{Mathematics}
    Jacquenetta preserves serif fonts in math mode for readability.

    \vspace{0.2cm}
    The posterior probability via Bayes' theorem:
    \[
        P(\theta \mid \mathcal{D}) =
        \frac{P(\mathcal{D} \mid \theta)\, P(\theta)}
             {P(\mathcal{D})}
    \]

    An inline example: given $n$ samples, the MLE is
    $\hat{\theta} = \arg\max_\theta \mathcal{L}(\theta; x_1, \dots, x_n)$.

    \vspace{0.2cm}
    Multi-line alignment:
    \begin{align}
        \mathcal{L}(\theta)
            &= \textstyle\sum_{i=1}^{n} \log p(x_i \mid \theta)
        &
        \nabla_\theta \mathcal{L}
            &= \textstyle\sum_{i=1}^{n}
                \nabla_\theta \log p(x_i \mid \theta)
    \end{align}
\end{frame}

% ============================================================
\section{Elements}
% ============================================================

\sectionframe{02}{Elements}

\begin{frame}{Lists}
    \begin{columns}[T]
        \begin{column}{0.45\textwidth}
            \textbf{Itemize}
            \begin{itemize}
                \item First level item
                \item Another item
                    \begin{itemize}
                        \item Second level
                        \item Nested item
                    \end{itemize}
                \item Back to first level
            \end{itemize}
        \end{column}

        \begin{column}{0.45\textwidth}
            \textbf{Enumerate}
            \begin{enumerate}
                \item First step
                \item Second step
                    \begin{enumerate}
                        \item Sub-step A
                        \item Sub-step B
                    \end{enumerate}
                \item Third step
            \end{enumerate}
        \end{column}
    \end{columns}

    \vspace{0.4cm}
    \textbf{Description}
    \begin{description}
        \item[Precision] Fraction of true positives among predicted positives.
        \item[Recall] Fraction of true positives among actual positives.
    \end{description}
\end{frame}

\begin{frame}{Blocks}
    All blocks share a unified left-border style, distinguished by color.

    \begin{block}{Default block}
        Used for general information, definitions, or neutral content.
    \end{block}

    \begin{alertblock}{Alert block}
        Used to highlight warnings, caveats, or important remarks.
    \end{alertblock}

    \begin{exampleblock}{Example block}
        Used for examples, demonstrations, or supporting evidence.
    \end{exampleblock}
\end{frame}

\begin{frame}{Custom environments}
    \begin{problock}{Problock --- callout with accent border}
        Ideal for theorems, definitions, or key takeaways.
        For any events $A$ and $B$ with $P(B) > 0$:\;
        $P(A \mid B) = {P(B \mid A)\, P(A)}/{P(B)}$
    \end{problock}

    \begin{highlightbox}
        \textbf{Highlightbox} --- an orange callout for observations,
        warnings, or results that need to stand out.
    \end{highlightbox}

    \vspace{0.2cm}
    Both use the same left-border language as standard blocks,
    but with tcolorbox for richer content (math, nested lists, etc.).
\end{frame}

\begin{frame}{Tables}
    Use \texttt{booktabs} for clean, professional tables.

    \vspace{0.4cm}
    \centering
    \begin{tabular}{lccc}
        \toprule
        \textbf{Model} & \textbf{Accuracy} & \textbf{F1 Score} &
        \textbf{Time (s)} \\
        \midrule
        Logistic Regression & 0.89 & 0.87 & 0.3 \\
        Random Forest       & 0.93 & 0.91 & 2.1 \\
        Neural Network      & 0.95 & 0.94 & 15.7 \\
        \textbf{Ensemble}   & \textbf{0.96} & \textbf{0.95} & 18.1 \\
        \bottomrule
    \end{tabular}
\end{frame}

\begin{frame}{Columns}
    \begin{columns}[T]
        \begin{column}{0.48\textwidth}
            \begin{problock}{Left column}
                Columns are useful for placing text and figures
                side by side, or for comparing two approaches.
            \end{problock}

            \vspace{0.2cm}
            Key advantages:
            \begin{itemize}
                \item Better use of slide space
                \item Visual separation of ideas
                \item Natural comparison layout
            \end{itemize}
        \end{column}

        \begin{column}{0.48\textwidth}
            \begin{problock}{Right column}
                You can mix any content: text, math, figures,
                tables, or code listings.
            \end{problock}

            \vspace{0.2cm}
            A quick formula:
            \[
                \text{ELBO} = \mathbb{E}_{q}[\log p(x \mid z)]
                - D_{\text{KL}}(q \| p)
            \]
        \end{column}
    \end{columns}
\end{frame}

% ============================================================
\section{Colors}
% ============================================================

\sectionframe{03}{Colors}

\begin{frame}{Color palette}
    Jacquenetta defines four theme colors:

    \vspace{0.4cm}
    \begin{columns}[T]
        \begin{column}{0.22\textwidth}
            \centering
            \begin{tikzpicture}
                \fill[jqdark] (0,0) rectangle (2.5,1.5);
                \node[white] at (1.25,0.75) {\textbf{Dark}};
            \end{tikzpicture}\\[4pt]
            {\small\texttt{\#121212}}
        \end{column}
        \begin{column}{0.22\textwidth}
            \centering
            \begin{tikzpicture}
                \fill[jqgray] (0,0) rectangle (2.5,1.5);
                \node[white] at (1.25,0.75) {\textbf{Gray}};
            \end{tikzpicture}\\[4pt]
            {\small\texttt{\#757575}}
        \end{column}
        \begin{column}{0.22\textwidth}
            \centering
            \begin{tikzpicture}
                \fill[jqblue] (0,0) rectangle (2.5,1.5);
                \node[white] at (1.25,0.75) {\textbf{Blue}};
            \end{tikzpicture}\\[4pt]
            {\small\texttt{\#4A90E2}}
        \end{column}
        \begin{column}{0.22\textwidth}
            \centering
            \begin{tikzpicture}
                \fill[jqorange] (0,0) rectangle (2.5,1.5);
                \node[white] at (1.25,0.75) {\textbf{Orange}};
            \end{tikzpicture}\\[4pt]
            {\small\texttt{\#E67E22}}
        \end{column}
    \end{columns}

    \vspace{0.6cm}
    Use them directly: \texttt{\textbackslash color\{jqblue\}} or
    \texttt{\textbackslash textcolor\{jqorange\}\{...\}}.

    \vspace{0.2cm}
    The \textbf{accent} color (default: {\color{jqaccent}blue})
    controls the \texttt{problock} border and standard block titles.
\end{frame}

% ============================================================
\section{Conclusion}
% ============================================================

\sectionframe{04}{Conclusion}

\begin{frame}{Summary}
    \begin{problock}{Jacquenetta at a glance}
        \begin{itemize}
            \item Minimalist design with a signature dark border
            \item Clean typography: Helvetica body, serif math
            \item Custom environments: \texttt{problock} and
                  \texttt{highlightbox}
            \item Section dividers with \texttt{\textbackslash sectionframe}
            \item Configurable accent color and optional border
        \end{itemize}
    \end{problock}

    \vspace{0.3cm}
    \begin{highlightbox}
        \centering
        \textbf{Get started:}\quad
        \texttt{\textbackslash usetheme\{Jacquenetta\}}
    \end{highlightbox}
\end{frame}

\thanksframe

\end{document}
